\section{The Gromov--Milman theorem}\label{sec:gromov-milman}

The first natural example of an extremely amenable group came from asymptotic
geometric analysis: concentration of measure on high-dimensional spheres and
groups forces fixed points.

\begin{theorem}[Gromov--Milman \cite{GromovMilman1983}]\label{thm:gromov-milman}
\uses{def:ea}
The unitary group $U(\ell^2)$ of an infinite-dimensional separable Hilbert
space, with the strong operator topology, is extremely amenable.
\end{theorem}

\begin{remark}[Lévy groups]\label{rem:levy}
\uses{thm:gromov-milman}
The proof shows that $U(\ell^2)$ is a \emph{Lévy group}: roughly, it contains
an increasing chain of compact subgroups with dense union whose Haar measures
concentrate, in the sense of Lévy's concentration of measure phenomenon.
In \cite{GromovMilman1983} the fixed point theorem is stated for increasing
unions of compact groups whose normalised Haar measures form a Lévy family,
acting equicontinuously on a compact space (Theorem~7.1 there), and it is
applied to the union of the groups $SO(n)$, to the group of finitary
permutations of $\mathbb{N}$ and to the union of the groups
$(\mathbb{Z}/2\mathbb{Z})^n$. Glasner observed that equicontinuity is not
needed \cite{Glasner1998}, so every Lévy group is extremely amenable.
Theorem~\ref{thm:gromov-milman} is the form in which the result is usually
quoted \cite{Uspenskij2000,GiordanoPestov2007}. This mechanism, which Pestov
calls the Ramsey--Dvoretzky--Milman phenomenon, gives many further examples
\cite{Pestov2006}.
\end{remark}
